% ===== §7 Praxis: cultivating a generative field =====
\section{Praxis: The Cultivation of a Generative Field}
\label{sec:praxis}

\noindent
The paper turns to practice, to the question latent in the preceding development: the manner of action within a system that reflexively reconfigures, and may alienate, the subject that acts within it. The ontology imposes a constraint that excludes the standard model of practice, and the determinate form of the practice, a cultivation rather than a construction, and a subtraction rather than an addition, is established only once that model is excluded.

\subsection{The exclusion of the external operator}
\label{sec:no-operator}

The standard model of practice is cybernetic: a subject external to the system observes its dynamics and intervenes to direct it toward the good. The model is unavailable, and its unavailability is a consequence of the preceding development rather than an incidental restriction. \xref{sec:reflexive} and \xref{sec:subject-in-cycle} established that there is no subject external to the cycle: the practitioner is internal to the dynamics, is recurrently reconfigured by it, and is a product of the system the practitioner would act upon. There is no completed self external to the process and operating it.

This excludes every practice of external mastery, and establishes that such a practice would in any case be self-defeating. To stand external to the relation and to administer it as an object to be optimised, extracted from, or engineered toward a specified output is the posture the geometry of \pinkref{Paper~IX} identified as the vicious cycle: the treatment of the relation as a possessible object. The posture of external control is itself the structure of the vicious cycle; the intervention from an external position to compel the good effects, by the form of its operation, the depletion it would prevent. The question is therefore reconstituted: not the regulation of the dynamics toward the good from an external position, there being no external position and the external posture being itself the depletion, but the manner of participation, within a reflexive system that may alienate the participant, such that the participation tends toward the good cycle. The regulator is the regulated; the practice is not a force applied to the dynamics but a tendency internal to it.

\subsection{Cultivation of a field in place of construction of a grammar}
\label{sec:field-not-grammar}

This reconstitution corrects a description that requires correction. One speaks loosely of constructing the shared grammar, of two parties making the rules of their relation. But the construction of a grammar is the legislation of productions, the antecedent fixing of the manner in which the relation shall generate, which is the settling and possession identified in \xref{sec:grammar-ontological}: the externalist's notation applied to the dynamics. The grammar is not constructed but cultivated into an openness within which it admits of continual joint reconfiguration. What the practice tends is not the content of the grammar, which is generated, jointly, by the parties within the field, but the condition under which the surface grammar admits of continual joint reconfiguration: the openness of the field, its curvature, the conditions under which the positive phase arises.

The term \emph{field} is employed with precision and not as loose figure. A field does not prescribe the trajectory of what moves within it; it determines the curvature, the potential, the conditions of motion, and admits the unfolding of each entity's dynamics along them. To cultivate a field rather than construct a grammar is to effect, for a relation, what a field effects for a particle: not the dictation of the path but the determination of the space within which paths are freely taken. The good the paper specifies is a condition of the space, a curvature under which the positive phase arises, and not a prescribed path. This is the practical form of the eudaimonic thesis of the prior papers: a good relation is not a structure two parties complete but a field within which each party's dynamics unfolds, flourishing being the unfolding of a subject's own dynamics within such a field. The practitioner does not author the content of the relation; the practitioner tends the field within which that content is jointly authored.

\subsection{The subtractive character of the practice}
\label{sec:subtractive}

The tending of such a field admits of determinate specification, and the actions specified share a form the preceding development establishes: they are negative, subtractive, refusals rather than productions, the positive phase being non-producible directly without its destruction (\pinkref{Paper~IX}), such that the available practice is the removal of what obstructs the curvature under which the positive phase arises. Three such tendings are specified.

The \emph{first} is the preservation of divergence. The criterion of a living grammar is its admission of genuine divergence, the capacity of the other party to take an unanticipated next step, to become an unanticipated subject, without penalty and without the withdrawal of regard as the cost of the divergence. The preservation of divergence is the refusal to narrow the relation's admissible productions to a single anticipated path, sustained as the recurrent question whether, within the relation, the other party retains the capacity to become an unanticipated subject. When one party's next move admits of a single permitted derivation, converging upon a prescription, the grammar has ceased to generate and generation has become regulation. This preservation is not only the cultivation of the field; it is, as \xref{sec:keystone} establishes, the practitioner's sole means of correcting the practitioner's own alienation, the other party's capacity to register extraction being the only available instrument by which the practitioner's own displacement is detected.

The \emph{second} is the sustaining of the opening without its occupation. This is the most counter-intuitive of the three: confronted with the unsayable term of the relation, the injury, the vulnerability, the present that resists articulation, the practice is to remain with it, to sustain it, and not to occupy it with explanation, designation, or premature commitment. The occupation of the opening, its settling with a sign, confers the security of determination and extinguishes the real value that subsists only around the opening (\xref{sec:opening}). The tending of the field is the sustaining of the opening, such that the positive phase is undergone recurrently rather than cashed once in a designation.

The \emph{third} is the preservation of the made language as poetic rather than as capital. Two parties generate private signs, a shared idiom, an internal designation, a metaphor legible only to them; these are joint attention condensed into a re-traversable poetic loop, the manifest trace of the grammar's joint reconfiguration, and the creation of language the paper's title denotes. The practice preserves these signs as poetic, sustaining the opening, soliciting the re-traversal, accumulating a fresh remainder at each traversal, rather than permitting their hardening into exclusionary symbolic capital, in which the legibility internal to the parties becomes a barrier against the exterior and a regulation within. The made language is well tended when it remains a solicitation to traverse the real, and ill tended when it sets into a settled token to be cashed, displayed, or deployed.

\subsection{An illustrative case}
\label{sec:concrete-field}

The three tendings are exhibited in a single situation, described in the manner of \xref{sec:phenomenology}.

\begin{casebox}{the tended field under strain}
Two parties, after some years, possess a private language: designations, gestures, a shared mode of falling silent each can read. A period of strain supervenes. The available recourse is control: the exhaustive designation of the difficulty, and thereby its mastery, which occupies the opening; the requirement that the other party cease to diverge and resume the anticipated role, which narrows the productions; the invocation of the shared idiom as a claim, the prior commitment, the established definition of the relation, which cashes the relation's accumulated tokens against the present difficulty, hardening the poetic into capital. Each is an addition, a compulsion, an operation from an external position; each, by the geometry, deepens the negative curvature it is intended to repair. The cultivation is the contrary, and is subtractive throughout: the difficulty is left partly unsaid rather than mastered by designation; the other party is permitted, in this period, to be unanticipated, without withdrawal of regard; the established idiom is re-traversed afresh rather than cashed as a claim. None of this produces the recovery. It removes the compulsions under which recovery cannot arise, and admits the field's curvature, the relation's own dynamics, to effect the positive phase no compulsion could.
\end{casebox}

\subsection{The internal criterion}
\label{sec:criterion}

The practice requires a means of determining, internally, whether the field tends toward the good or its contrary, a criterion held internally, posed as a question, and not a metric read from an external position, the latter restoring the excluded external operator. Two questions are sufficient, and they are the orthogonal pair the series has developed: the poetic criterion and the criterion of justice, neither reducible to the other.

The first is the dynamical, poetic question: whether the loop continues to spiral or merely recurs. When the private signs, the practices, the shared silences are re-traversed, whether they continue to accumulate an irreducible remainder, the spiral and positive holonomy, or have emptied into a performance recurring for its own sake, the circle and zero holonomy. The second is the question of justice: whether the spiral's positive phase excludes any party from its benefit, present solely as its cost. Returning to the political economy of \xref{sec:political}: whether, in the coupled system of the two parties and the families behind them, there is a party, characteristically the party bearing the labour of care and reproduction, idealised as natural devotion, that creates value and is never reached, positively, by its return; that is reconfigured solely negatively; that is, within the system, solely fuel.

\begin{claim}[The practice as the posing of the two questions and subsequent subtraction]
The practice is the recurrent and exacting posing of the two orthogonal questions, whether the loop continues to spiral, and whether any party is solely its cost, and, when the answers turn negative, the operation of subtraction: the withdrawal of the extraction, the relaxation of the compelled closure, the restitution of the possessed phase. The practice does not add, optimise, or produce; it removes what obstructs the curvature under which the good arises. This is the content of \emph{wu wei} for two parties: not the absence of practice but the redirection of practice from the production of the good relation to the non-obstruction of the field within which a good relation generates, and from a control exercised over the dynamics to a tending exercised within it.
\end{claim}

\noindent
A constraint remains, which the section has twice indicated and which is now stated. The practice requires of the practitioner the determination of whether the field tends toward the good, the subtraction of the extraction, the preservation of the other's divergence. But the practitioner is a product of the cycle, and possibly a product of an alienated cycle, the developed form of which is precisely that it is not registered as alienation, the extraction being experienced as regard. A subject already reconfigured by a vicious cycle poses the two questions in good faith and obtains reassuring answers, taking itself to be tending the field while depleting it. The entitlement of such a practitioner to the practice of the good is thereby in question. The practice, posed exactingly, encounters a limit it cannot cross from within itself, and that limit is the matter of the following section.
